\begin{center}
\begin{tikzpicture}[x=.88cm,y=.88cm,font=\small]
  \fill[paper] (0,0) rectangle (12.2,6.35);

  \fill[dynasty!24] (.8,1.05) -- (3.25,.8) -- (4.0,2.2) -- (3.15,3.45) --
    (1.05,3.2) -- cycle;
  \node[align=center] at (2.35,2.18) {关中\\秦旧行政与粮仓};

  \fill[timeline!27] (1.25,4.05) ellipse (1.25 and .73);
  \node[align=center] at (1.25,4.05) {汉中\\刘邦受封};

  \fill[source!22] (4.45,2.25) ellipse (1.12 and .9);
  \node[align=center] at (4.45,2.25) {三秦\\雍、塞、翟};

  \fill[dynasty!42] (8.15,3.9) ellipse (1.55 and 1.05);
  \node[align=center,text=white] at (8.15,3.9) {西楚\\项羽的中心};

  \fill[timeline!24] (10.65,1.65) ellipse (1.12 and .78);
  \node[align=center] at (10.65,1.65) {诸侯集团\\各自兵众};

  \draw[-{Stealth[length=2mm]},very thick,color=dynasty]
    (2.15,3.48) .. controls (2.45,3.75) and (2.95,3.7) .. (3.45,3.0);
  \node[text=dynasty,align=center,font=\scriptsize] at (3.12,4.08) {由汉中\\出兵还定};

  \draw[-{Stealth[length=2mm]},very thick,color=dynasty]
    (5.6,2.62) .. controls (6.25,3.42) and (6.7,3.65) .. (7.05,3.75);
  \node[text=dynasty,above,font=\scriptsize] at (6.28,3.7) {竞争关中};

  \draw[-{Stealth[length=2mm]},thick,color=timeline]
    (9.52,3.27) .. controls (9.95,2.75) and (10.22,2.4) .. (10.45,2.25);
  \draw[-{Stealth[length=2mm]},thick,color=timeline]
    (9.52,4.47) .. controls (10.05,4.65) and (10.5,4.37) .. (10.98,3.88);
  \node[text=timeline,align=center,font=\scriptsize] at (10.85,4.88) {分封与\\联盟牵制};

  \node[draw=source!75,fill=paper,rounded corners=2pt,align=center,
    inner sep=3pt,font=\scriptsize] at (6.18,.85)
    {前206年的分封并未固定胜负\\而是把关中资源与诸侯关系重新拆开};
\end{tikzpicture}

\smallskip
\small\textit{图\quad 前206年前后的竞争中心示意：据《史记·项羽本纪》《高祖本纪》及《汉书·高帝纪》《异姓诸侯王表》，概括项羽分封、刘邦受封汉中与还定三秦后形成的几个动员中心。节点仅示相对关系，非等比例地图；三秦、汉中、西楚及诸侯集团的轮廓不表示精确疆界，箭线亦不重建具体行军路线。}
\end{center}
